% ============================================================
% 字体配置文件 (macOS)
% 使用 XeLaTeX 引擎，通过 fontspec 设置中英文字体
% ============================================================
\usepackage{fontspec}
\usepackage{xeCJK}

% ---------- 英文字体（衬线 / 无衬线 / 等宽） ----------
\setmainfont{Times New Roman}
\setsansfont{Arial}
\setmonofont{Courier New}

% ---------- 中文字体 ----------
% ===== 中文正文（宋体）=====
\setCJKmainfont{Songti SC}[
  BoldFont       = Heiti SC,
  ItalicFont     = Kaiti SC,
  BoldItalicFont = Heiti SC
]

% ===== 中文无衬线（黑体，标题）=====
\setCJKsansfont{Heiti SC}

% ===== 中文等宽 =====
\setCJKmonofont{STFangsong}

% ---------- 数学字体 ----------
% 如需自定义数学字体，需先在 packages.tex 中加载 unicode-math 宏包，
% 然后取消下面的注释：
% \setmathfont{Latin Modern Math}
